\subsection{Назначение и среда работы}

Программа предназначена для автоматической двусторонней синхронизации файлов между двумя каталогами на разных компьютерах. На каждом компьютере запускается экземпляр программы. Экземпляры взаимодействуют непосредственно по сети с помощью TCP-сокетов; выделенный сервер и СУБД не требуются. Для запуска необходима Java версии 17 или новее.

\subsection{Исходные данные и управление}

Пользователь указывает существующий локальный каталог, адрес второго компьютера, локальный TCP-порт, порт второго экземпляра и интервал проверки в секундах. Настройки вводятся в графическом окне либо передаются параметрами командной строки. После запуска каждый экземпляр принимает входящие соединения и периодически инициирует обмен со вторым экземпляром. Доступна команда немедленной проверки. При недоступности узла следующая попытка выполняется через заданный интервал.

\subsection{Правила обработки файлов}

Программа рекурсивно формирует список обычных файлов в каталоге. Для каждого файла сохраняются его относительный путь, размер и хэш SHA-256. Экземпляры обмениваются списками и сравнивают файлы по относительным путям.

\begin{itemize}
  \item Файл, присутствующий только в одном каталоге, передаётся в другой каталог с сохранением относительного пути.
  \item Файлы с одинаковым путём и одинаковым содержимым повторно не копируются.
  \item При одинаковом пути и различном содержимом возникает конфликт: обе версии сохраняются, а сообщение с путём и хэшами записывается в журнал.
  \item Получаемые данные сначала записываются во временный файл и проверяются по размеру и хэшу. Существующий файл не перезаписывается.
\end{itemize}

Удаления между каталогами не распространяются: отсутствующий файл при следующей проверке копируется снова. Пустые каталоги отдельно не передаются. Символические ссылки не синхронизируются. Каталог \texttt{.filesynchronizer-tmp} в корне выбранного каталога зарезервирован для временных файлов и исключён из обмена; собственный файл журнала также исключается. Имена с обратной косой чертой или двоеточием не поддерживаются, поскольку такие имена несовместимы с файловыми системами разных ОС.

\subsection{Журнал и сообщения об ошибках}

В текстовый журнал записываются запуск и остановка сервера, адреса подключившихся клиентов, сведения о переданных файлах, результаты проверок, конфликты и ошибки. В графическом режиме эти сообщения одновременно отображаются в окне. Ошибки настройки выводятся пользователю. При разрыве соединения незавершённый файл не устанавливается как результат, а обмен повторяется при следующей проверке.

\subsection{Ограничения и требования к реализации}

Соединение не шифруется и не проверяет личность удалённой стороны, поэтому программа рассчитана на доверенную сеть. Реализация разделена на модель, представление, контроллер, сетевую подсистему и вспомогательные средства. Приложение и тесты используют стандартную библиотеку Java. Проект собирается Maven; сборка компилирует программу, выполняет тесты модели и сетевого взаимодействия и создаёт JAR-файлы программы и тестов.
